\section{Introdução}
\label{sec:introducao}

% definicao e ref que baseou o blockchain

\PARstart{O}{problema} que o \textit{Blockchain} busca resolver é a necessidade de um sistema distribuído confiável, onde não é possível alterar o conteúdo e registros de data e hora de um conjunto de dados sem que seja perceptível a alteração \cite{di2017blockchain}. Essa tecnologia de encadeamento de blocos foi primeiramente apresentada para suporte ao \textit{Bitcoin}, o qual foi criado para ser uma forma alternativa às instituições financeiras para processamento de pagamentos eletrônicos, utilizando um sistema de assinuturas digitais conectados à uma rede \textit{peer-to-peer} \cite{nakamoto2008bitcoin,rocha2004peer}. 

% contextualizacao do problema

Para uma melhor compreensão da aplicabilidade do \textit{Blockchain}, utilizando de um exemplo do mercado imobiliário do artigo de Di Pierro \cite{di2017blockchain}, imagine o caso de vendas de casas, onde o dono é definido por aquele em que a casa foi vendida pela última vez, e a propriedade só pode ser verificada por um conjunto de registros que comprovem todas as transações relacionadas à casa (livro-razão), o qual geralmente é mantido em um cartório de imóveis. O problema que surge é a dependência exclusiva sob uma única instituição para manter esse livro-razão seguro de ações fraudulentas. Então, para solucionar isso, o \textit{Blockchain} fornece um mecanismo distribuído com múltiplos \textit{peers} que mantêm uma cópia do livro-razão e são capazes de verificar se um registro e sua respectiva datação sofreram adulteração \cite{di2017blockchain}.

% diversificação do blockchain e aplicações

Apesar da aplicação mais conhecida ser o \textit{Bitcoin}, a técnica de encadeamento foi utilizada em vários serviços financeiros, como ativos digitais, remessas e pagamentos \textit{online} \cite{trendsCrypto}. Além do mercado financeiro, o crescimento do \textit{Blockchain} ocorreu devido à sua capacidade de ser útil em diversas áreas da indústria, como a da saúde, onde a partir da digitalização de registros médicos surgiu a necessidade de que os dados adquiridos em diferentes ambientes hospitalares fossem compartilhados com segurança \cite{blockSurvey1,blockMed1}. Neste trabalho, o estudo de \textit{Blockchain} foca no domínio financeiro, devido ao seu pioneirismo, e no domínio médico, em alinhamento com a área de pesquisa acadêmica atual do autor. Apesar desse recorte de escopo, ressalta-se que todos os seguintes domínios foram influenciados por essa tecnologia: logístico \cite{Hackius2017}, manufatura \cite{Barenji2018}, energia \cite{Mengelkamp2018}, agricultura \cite{Lin2017}, robótica \cite{Ferrer2018}, entretenimento \cite{Kalra2018}, construção \cite{Wang2017} e telecomunicação \cite{Ibrahim2017}.

% impacto no universo de sistemas distribuídos para finalizar o contexto

A partir desse cenário de expansão do \textit{Blockchain}, nota-se uma mudança no universo de sistemas distribuídos, que inicialmente os algoritmos de consenso eram projetados para ambientes computacionais confiáveis, fechados e sob o controle de organizações específicas, onde a principal preocupação eram falhas técnicas comuns ou interrupções de hardware \cite{Li2021}. O \textit{Blockchain} redefiniu esse cenário ao transferir o desafio do consenso para redes abertas, em grande escala e desprovidas de confiança mútua, mantendo a integridade dos registros sem a dependência exclusiva de autoridades centrais \cite{Xiao2020}.
